\begin{proof}
\begin{enumerate}[label=\arabic*.]
\item The critical ideal of $h$ is $\mathcal I_D$. Locally
$q=x_0\cdots x_r$, its generators are the products omitting
one factor. A vector field in the kernel of $dq$ has coefficients
$a_i=x_i b_i$ in the crossing directions and
$\sum b_i=0$; the remaining ordinary directions are free.
Thus the kernel is the translation bundle $h^*E_0$, and
\[
0\to h^*E_0\to T_M\to J\to0,
\qquad J=\mathcal I_D\otimes h^*L^9.
\]
The conductor square gives
\[
0\to\mathcal O_W\to\nu_*\mathcal O_Y\oplus\mathcal O_D
 \to\nu_*\mathcal O_{\partial Y}\to0.
\]
The boundary is connected and anticanonical in the toric
threefold $Y$. Its coherent cohomology is $(1,0,1)$, by
$0\to K_Y\to\mathcal O_Y\to\mathcal O_{\partial Y}\to0$
and Serre duality. The singular locus $D$ is connected and
has dimension two, as is immediate from the connected
double-face incidence graph. The map on constants is
$(a,b)\mapsto a-b$ and is onto. In higher degree the conductor
sequence makes $H^1(\mathcal O_W)\to H^1(\mathcal O_D)$
an isomorphism. The map
$H^2(\mathcal O_{\partial Y})\to H^3(\mathcal O_W)$ is onto,
since $H^3(\mathcal O_D)=H^3(\mathcal O_Y)=0$; both sides
have dimension one, so it is an isomorphism. Consequently
$H^2(\mathcal O_W)\to H^2(\mathcal O_D)$ is also an
isomorphism and $h^q(\mathcal O_D)=(1,3,3)$.

\item Projection and \eqref{eq:P1Adegree} give
\[
h^q(h^*E_0)=(1,1,0,0,0),\qquad
h^q(h^*L^9)=(1,2,1,0,0).
\]
The maps from the latter groups to $H^q(\mathcal O_D)$
are evaluation at the cusp followed by the just-proved
restriction isomorphisms. Each contributing coarse summand
has degree zero, so these evaluations are injective, with
ranks $1,2,1$. The critical-ideal sequence gives
$h^q(J)=(0,0,1,2,0)$. The tangent dimensions and the
isomorphism $H^1(h^*E_0)\to H^1(T_M)$ follow.

\item Differentiating the actual periods gives, on the universal
torus covering, the Cech cocycle of constant translations
obtained by differentiating each lattice period. Under the
constant-cocycle description of fibre cohomology this is
the map $F\to V/F$ sending $f$ to $\alpha(f)e$.
It is nonzero and occupies precisely the trivial summand
of $E_0\otimes G$. The Leray edge map identifies this summand
with $H^1(h^*E_0)$, so the global Kodaira--Spencer class is
nonzero. The finite and toroidal charts depend holomorphically
on the parameter and supply the actual compact family,
so this cocycle is the derivative of that family on every chart.

\item Douady's analytic versality and semiuniversality theorems
\cite[\S VII.8, pp.~598--599; \S VIII.1, pp.~599--600]{Douady1974} apply to the
compact complex manifold. The Kuranishi germ embeds in its
one-dimensional tangent space. A classifying map from the
actual parameter disc has nonzero derivative, hence its
composite with that embedding is an invertible one-variable
analytic germ. Its pullback on $\mathbb C\{t\}$ is injective,
so the defining ideal of the Kuranishi germ is zero, including
nilpotents. This proves smoothness and miniversality.

\end{enumerate}
\end{proof}
